% Source: https://artofproblemsolving.com/community/c3h2565913_a3b3a2b2
% Stable identity: aops-topic-2565913-q01
If $a=0$, then $b^3=0$, so $b=0$; similarly $b=0$ forces $a=0$. This gives $(0,0)$.

Now let $a,b\ne0$, let $d=\gcd(a,b)>0$ and write $a=dx$, $b=dy$ with $\gcd(x,y)=1$. Dividing $d^3\left(x^3+y^3\right)=d^4x^2y^2$ by $d^3$,
\[x^3+y^3=d\,x^2y^2.\]
Here $x^2$ divides $x^3$ and $dx^2y^2$, so $x^2\mid y^3$. As $\gcd(x,y)=1$, also $\gcd\left(x^2,y^3\right)=1$, so $x^2=1$. Likewise $y^2=1$. Then $x^3+y^3=d>0$, which forces $x=y=1$ and $d=2$, i.e. $a=b=2$. Indeed $8+8=16=4\cdot4$.

The solutions are $\boxed{(a,b)=(0,0)\text{ and }(2,2)}$.
